\problemname{Protein Synthesis}
It is a well-known fact that Latbo is the laziest city in Sweden. Ever since he was a child, Björn has dreamed
of becoming the laziest person in town. Yesterday, he made it all the way to the final of this year's
laziness competition. He was fully prepared to win, but was utterly crushed when his opponent was
too lazy to even show up.

This was the last straw for Björn. He has now decided to use his medical knowledge to develop a cure
for laziness. To do this, he needs to synthesize the magical protein $T$. However, old habits die hard
as he tries to synthesize it using as few other proteins as possible.

Since he's too lazy to figure out how to do this himself, he now wants your help to find the best way to
construct protein $T$. Protein $T$ can be represented as a string of length $N$. Björn knows how to
produce $M$ types of proteins and wants to minimize the number of proteins he needs to manufacture.
Once the correct set of proteins has been produced, it's easy to assemble them into protein $T$.

\section*{Input}
The first line of the input contains the integer $N$ ($1 \le N \le 10^5$), the length of protein $T$.

The next line contains a string of length $N$, representing protein $T$.

The next line contains $M$ ($1 \le M \le 10^6$), the number of proteins Björn can manufacture.

Afterwards, $M$ lines follow, each containing a string. This string represents a protein he can produce.
It is guaranteed that the total length of the $M$ strings does not exceed $10^6$.

All strings in the input consist of letters \texttt{a}-\texttt{z}.

\section*{Output}
Print an integer: the minimum number of proteins Björn needs to produce to create protein $T$.

\section*{Points}
Your solution will be tested on several test case groups.
To get the points for a group, it must pass all the test cases in the group.

\noindent
\begin{tabular}{| l | l | p{12cm} |}
  \hline
  \textbf{Group} & \textbf{Points} & \textbf{Constraints} \\ \hline
  $1$    & $10$       & $N \leq 20, M \leq 5$ \\ \hline
  $2$    & $20$       & $N, M \leq 100$ \\ \hline
  $3$    & $20$       & $M \leq 100$ \\ \hline
  $4$    & $40$       & $N \leq 10,000$ \\ \hline
  $5$    & $10$       & No additional constraints. \\ \hline
\end{tabular}

\section*{Explanation of sample 1}
By creating the protein \texttt{abc} twice, he can combine them to create protein $T$.
